\chapter{Introdução}
\label{cap:Introdução}

A análise de dados permite o conhecimento, por exemplo, de tendências, gargalos e predição de problemas, entre outras informações úteis, auxiliando na tomada de decisões estratégicas em vários setores da economia.  O acesso a essas informações por parte de empresas e governos permite que tais decisões sejam melhor embasadas, reduzindo custos, melhorando a distribuição de mercadorias e serviços, e podendo, inclusive, até salvar vidas.

Atualmente, é comum que organizações possuam uma miríade de ativos de TI (Tecnologia da Informação) à disposição, como computadores, impressoras, equipamentos de rede e \textit{softwares}, por exemplo, utilizados pelos trabalhadores. Fazer a gestão de todos esses ativos é um problema que muitas empresas enfrentam ao longo de suas vidas. Entre outros pontos, é necessário saber quais são os usuários dos recursos, se determinado recurso continua ou não sendo útil, se seu prazo de vida útil já se esgotou, se algum equipamento está apresentando mais defeitos que a média, além de controlar possíveis alterações feitas pelo setor de TI no ambiente de usuário.  Para auxiliar nessa tarefa, utilizam-se ferramentas que permitem a análise e visualização de dados junto aos sistemas ITSM (\textit{Information Technology Service Management} - Serviço de Gerenciamento de Tecnologia da Informação), que permitem o controle de chamados e ativos pertencentes a uma organização. 

Com acesso total à fonte de dados dos sistemas ITSM, é possível alimentar ferramentas de análise, utilizando rotinas \textcolor{red}{(minha dúvida é em relação à concordância, porque, como está, é como se ETL fosse uma rotina específica, e não um padrão )} de tratamento de dados, como a ETL (\textit{Extract}, \textit{Load}, \textit{Transform} - Extrair, Transformar, Carregar), por exemplo, que permitem a manipulação e o tratamento do que está contido na base de dados. Após a obtenção e o tratamento, visualizações podem ser definidas para permitir que o gestor possa antecipar decisões e negociar com outros departamentos e organizações com os quais tenha contato.

Dessa forma, passa-se de um modelo de gestão reativo para uma abordagem proativa, otimizando o gerenciamento e reduzindo custos. Em vez de aguardar a falha estrutural de um equipamento para solicitar sua substituição ou manter licenças de \textit{software} ociosas, gerando despesas invisíveis, a gestão de TI passa a atuar fundamentada em evidências claras.

\textcolor{red}{Precisei reorganizar os 2 parágrafos seguintes porque as frases que apontam o que o trabalho apresenta estavam espalhadas. Dê uma olhada se ficou algo redundante.}

Nas seções seguintes, este relatório  detalha o trabalho realizado pelo autor durante seu estágio, com foco em Análise de Dados, na empresa Électricité de France (EDF), de novembro de 2022 a novembro de 2024. 
Nesse período, o autor atuou no setor de TI da Usina Termelétrica Norte Fluminense (UTENF), pertencente à EDF Brasil, onde teve contato com as áreas de Infraestrutura e Sistemas, tendo ficado um ano em cada uma. Na área de Infraestrutura, aprendeu sobre o funcionamento dos equipamentos de TI, bem como de TO (Tecnologia de Operação), participou da manutenção de sistemas críticos e realizou catalogação e manutenção de ativos da área. No período em Sistemas, aplicou os conhecimentos de Banco de Dados e de Estatística para desenvolver relatórios, fornecer acesso controlado às informações dentro do ambiente Microsoft e criar consultas, provendo \textit{views} específicas segundo a demanda das outras áreas.

É demonstrado neste trabalho como a modelagem dos indicadores-chave de desempenho (KPIs - \textit{Key Performance Indicators}) e a construção de painéis interativos transformam um vasto volume de dados brutos em inteligência acionável, garantindo a otimização contínua do ambiente tecnológico e o direcionamento assertivo dos investimentos da área. 
É apresentado também o processo de elaboração de um relatório para acompanhamento de chamados registrados no sistema ITSM da empresa e o processo de desenvolvimento utilizado pelo autor para definir os tratamentos de dados e a escolha de apresentação dos mesmos. São demonstradas as rotinas de ETL utilizando a ferramenta Power BI da Microsoft, bem como a confecção de relatórios para a obtenção de \textit{insights} e a previsão de tendências.

O trabalho está organizado da seguinte forma. 
Nesta seção são apresentados alguns detalhes sobre a empresa: sua atuação, histórico, missão e valores, além de informações sobre o estágio do autor. Na Seção 2 são introduzidos o ambiente de dados e as rotinas para tratamento e carregamento desses dados na plataforma Power BI. A Seção 3 detalha a confecção de um relatório para o departamento de TI, demonstrando a aplicação dos conceitos abordados, os elementos utilizados que são disponibilizados pela ferramenta Power BI, e os recursos de interação disponíveis. Por fim, na Seção 4, são apresentadas as conclusões.


\section{A EDF}


A Électricité de France, mais conhecida pela sigla EDF, é uma empresa estatal francesa do ramo de energia elétrica. Criada por decisão governamental em 8 de abril de 1946, a empresa  atualmente possui o monopólio de toda a grade de energia francesa sob seu ``guarda-chuva''. \cite{EDF_Historico}


A expansão da EDF se deu por investimentos do Plano Marshall e do governo francês durante o período conhecido como os ``Trinta Gloriosos'', o que a tornou global. Seu primeiro investimento em solo estrangeiro foi realizado em 1992. \cite{EDF_Historico}


Em 2004, o mercado de energia francês foi desregulamentado, fazendo com que parte da EDF fosse  disponibilizada na bolsa de valores durante o ano de 2005. Isso levou a empresa a adotar uma divisão societária de capital misto, mantendo o governo como acionista majoritário. \cite{EDF_Historico} 


Em 2022, a EDF foi reestatizada \cite{EDF_Reestatizacao}. Segundo o governo francês, ``a crise climática e a situação geopolítica requerem decisões firmes para garantir a independência da França e a soberania energética'' \cite{Sen_Frances}.


A Figura 1 apresenta a logo da EDF em 2025.
\begin{figure}[!ht]
	\centering
	\caption{Logo da EDF - 2025.}
	\includegraphics[width=0.25\linewidth]{figuras/edf_logo.jpg}
	\caption*{Fonte: \cite{Logo_EDF}.}
	\label{fig:placeholder Figura 1}
\end{figure}


\subsection{A EDF no Brasil}

A EDF atua no Brasil desde 2004 nos ramos de geração e de transmissão de energia.
Considerando a geração de energia, a empresa é dona de um vasto portfólio de ativos, possuindo 2,7GW de capacidade instalada, entre ativos solares, eólicos, hídricos e termelétricas. 
Esses ativos, que antes eram divididos entre as entidades EDF Brasil e EDF Renováveis, foram unificados em 2025, após um processo de reestruturação, formando a EDF Power Solutions \cite{EDF_PS_Unificacao,EDF_PS_Brasil}.
Já no setor de transmissão de energia, a EDF arrematou, por meio de sua então subsidiária UTE Norte Fluminense SA, 6km de linhas de transmissão, com capacidade de transmissão de 300MVA, servindo ao Complexo Porto do Açu e à Santo Amaro, ambos no estado do RJ \cite{Leilao_Transmissao}.


\subsection{Missão da EDF}
A missão da EDF é:

``Our raison d’être is to build a net zero energy future with electricity and innovative solutions and services, to help save the planet and drive wellbeing and economic development.''  \cite{EDF_PS_Missao}


Em tradução livre: 
Nossa razão de ser é construir um futuro com energia sem emissões de $\text{CO}_2$, utilizando eletricidade, soluções e serviços inovadores para ajudar a salvar o planeta e impulsionar o bem-estar e o desenvolvimento econômico.


Considerando sua missão, os valores da EDF se baseiam fortemente na promoção da transição energética, da inclusão, da diversidade social e de um ambiente de trabalho seguro e saudável.

\section{Sobre o estágio de TI}


Após a entrada na EDF, o primeiro dia do estagiário é dedicado ao \textit{onboarding} do time de Recursos Humanos (RH), com  reuniões introdutórias e vídeos instrucionais sobre as políticas da empresa. Caso o estagiário seja alocado em uma planta industrial, também é realizado um \textit{onboarding} de segurança junto ao time de QSMA (Qualidade, Meio Ambiente e Segurança). 


Em seguida, o estagiário é direcionado para um dos times de atuação dentro do departamento de TI, passando a integrar a equipe de Infraestrutura ou a de Sistemas.
O time de Infraestrutura tem como responsabilidades:  manter a rede da empresa segura e operacional, zelar pelos ativos de TI (sejam eles físicos ou licenças de \textit{softwares}), configurar e manter os \textit{data centers} e a infraestrutura de nuvem, disponibilizar Máquinas Virtuais com  serviços aos usuários, e prestar assistência ao time de Operações em assuntos que envolvam sistemas de TO. Já o time de Sistemas tem como funções: manter os \textit{softwares} corporativos operacionais, zelar pela integridade dos dados dos Bancos de Dados, realizar implementações e modificações dos processos dos departamentos mediante requisição, disponibilizar dados para criação de visualizações, executar modificações e manutenções nas visualizações dos departamentos, e controlar o acesso dos usuários às visualizações, garantindo que não haja acessos não autorizados.


Durante a atuação do autor, limitada à duração máxima do programa de 2 anos, o time de Sistemas era composto por um Especialista em Sistemas, dois Analistas terceirizados especializados no Sistema IBM MAXIMO, um DBA (\textit{Database Administrator}) terceirizado e um estagiário. Esse time respondia pelo gerenciamento dos sistemas responsáveis pelas operações da empresa e pelo controle de acesso aos dados fornecidos. 
Eram esperadas de seu estagiário: 
\begin{itemize}
	\item a criação  e a manutenção de relatórios para diversos departamentos;
	\item  a criação e a manutenção do \textit{gateway} de dados utilizado para acesso aos dados da empresa pelo serviço na nuvem do Power BI;
	\item a confecção de consultas SQL (com acessos controlados pelo Especialista)  para a disponibilização de fluxos de dados na ferramenta Power BI; e
	\item a criação de novos conectores no Power BI para permitir a obtenção de dados de fontes diversas.
\end{itemize} 